\section{关于本书}

本笔记依据 M.\ Spivak 的《流形上的微积分》(Calculus on Manifolds) 整理而成，内容以原书章节为纲，辅以必要的补充与排版整理，供复习与查阅之用.

\section{符号说明}

本节列出全书通用的记号与约定.

\begin{itemize}
    \item 零向量$(0,\cdots,0)$通常简记为$\mathbf{0}$.
    \item $\mathbb{R}^n$的通常基底是$e_1,\cdots,e_n$，其中$e_i=(0,\cdots,1,\cdots,0)$在第$i$个位置上是$1$.
    \item 如$T:\mathbb{R}^n\to\mathbb{R}^m$是一个线性变换，$T$关于$\mathbb{R}^n$与$\mathbb{R}^m$的通常基底的矩阵是$m\times n$矩阵$A=(a_{ij})$，其中$T(e_i)=\sum_{j=1}^m a_{ji}e_j$——$T(e_i)$的系数出现在矩阵的第$i$列.
    \item 如$S:\mathbb{R}^m\to\mathbb{R}^p$有$p\times m$矩阵$B$，则$S\circ T$有$p\times n$矩阵$BA$[这里$S\circ T(\mathbf{x})=S(T(\mathbf{x}))$].绝大多数线性代数书籍把$S\circ T$简记为$ST$.
    \item 如$\mathbf{x}\in\mathbb{R}^n$与$\mathbf{y}\in\mathbb{R}^m$，则$(\mathbf{x},\mathbf{y})$表示
    $$(x_1,\cdots,x_n,y_1,\cdots,y_m)\in\mathbb{R}^{n+m}.$$
\end{itemize}

为找出$T(\mathbf{x})$，我们来计算$m\times 1$矩阵
$$\begin{pmatrix}
    y_1\\
    \vdots\\
    y_m
\end{pmatrix}
=
\begin{pmatrix}
    a_{11},\cdots,a_{1n}\\
    \vdots\quad\ddots\quad\vdots\\
    a_{m1},\cdots,a_{mn}
\end{pmatrix}
\cdot
\begin{pmatrix}
    x_1\\
    \vdots\\
    x_n
\end{pmatrix},$$
则$T(\mathbf{x})=(y_1,\cdots,y_m)$.
